\chapter{Całki niewłaściwe}
\label{ma1:ch:improper-integrals}

Całkę Riemanna skonstruowaliśmy dla funkcji ograniczonych na przedziałach domkniętych o skończonej długości. Teraz zbadamy, kiedy granice takich całek pozwalają rozszerzyć konstrukcję na przedziały nieograniczone lub funkcje nieograniczone w pobliżu końców. Samo znalezienie funkcji pierwotnej nie rozstrzyga jeszcze zbieżności całki niewłaściwej. W dowodach dla jednego niewłaściwego prawego końca przyjmujemy $a\in\R$, $a<b\le\infty$. Zapis $r\to b$ od wnętrza oznacza $r\to b^-$ dla skończonego $b$ oraz $r\to+\infty$ dla $b=\infty$.

\section{Definicja i rozdzielanie osobliwości}
\begin{definicja}
\label{ma1:def:improper-infinite}
Niech $f:[a,\infty)\to\R$ będzie całkowalna w sensie Riemanna na każdym $[a,R]$, $R>a$. Całkę niewłaściwą
\[
 \int_a^\infty f(x)\,\dd x
 :=\lim_{R\to\infty}\int_a^R f(x)\,\dd x
\]
nazywamy zbieżną, jeżeli granica istnieje i jest skończona. W przeciwnym razie jest rozbieżna; może być rozbieżna do $+\infty$ albo do $-\infty$. Analogicznie dla funkcji całkowalnej na każdym $[R,b]$ definiujemy $\int_{-\infty}^b f(x)\,\dd x=\lim_{R\to-\infty}\int_R^b f(x)\,\dd x$.
\end{definicja}
\begin{definicja}
\label{ma1:def:improper-finite}
Niech $a<b$ i niech $f:[a,b)\to\R$ będzie całkowalna na każdym $[a,r]$, $a<r<b$. Całkę niewłaściwą definiujemy jako
\[
 \int_a^b f(x)\,\dd x:=\lim_{r\to b^-}\int_a^r f(x)\,\dd x,
\]
jeśli granica istnieje i jest skończona; w przeciwnym razie całka jest rozbieżna. Dla funkcji określonej na $(a,b]$ i całkowalnej na każdym $[r,b]$, $a<r<b$, stosujemy analogicznie granicę $\lim_{r\to a^+}\int_r^b f(x)\,\dd x$.
\end{definicja}
Nie zakładamy rzeczywistej nieograniczoności funkcji. Jeżeli $f$ ma całkowalne przedłużenie na $[a,b]$ i $|f|\le B$ na tym przedziale, to z addytywności i oszacowania całki (twierdzenia \ref{ma1:thm:riemann-additivity} i \ref{ma1:thm:riemann-order})
\[
 \left|\int_a^b f\,\dd x-\int_a^r f\,\dd x\right|
 \le B(b-r)\longrightarrow0.
\]
Całka niewłaściwa zgadza się ze zwykłą całką Riemanna przedłużenia. Wartość przypisana końcowi nie ma znaczenia (wniosek \ref{ma1:cor:riemann-finite-changes}).

\begin{definicja}
\label{ma1:def:improper-splitting}
Gdy oba końce przedziału $(a,b)$ są niewłaściwe (dopuszczamy $a=-\infty$, $b=\infty$), zakładamy całkowalność na wszystkich domkniętych podprzedziałach i wybieramy $c\in(a,b)$. Całka $\int_a^b f$ jest zbieżna, jeżeli obie całki $\int_a^c f$ i $\int_c^b f$ są zbieżne. Definiujemy ją jako ich sumę. Przy skończonej liczbie osobliwości we wnętrzu $a<c_1<\cdots<c_m<b$ zakładamy całkowalność na każdym domkniętym odcinku nieobejmującym osobliwości (z dopuszczeniem regularnych końców). Rozdzielamy przedział w każdym $c_j$, a każdy odcinek o dwóch niewłaściwych końcach dodatkowo w dowolnym regularnym punkcie wewnętrznym. Wymagamy zbieżności wszystkich otrzymanych całek jednostronnych osobno; wartość całki jest wtedy ich skończoną sumą.
\end{definicja}
\begin{proof}[Niezależność od punktu rozdzielenia]
Niech $c<d$ będą punktami regularnymi w tym samym przedziale bez osobliwości. Dla $a<r<c<d<s<b$ addytywność daje
\[
 \int_r^d f=\int_r^c f+\int_c^d f,\qquad
 \int_d^s f=\int_c^s f-\int_c^d f.
\]
Całka $\int_c^d f$ jest właściwa. Przechodząc niezależnie do granic $r\to a$ i $s\to b$ od wnętrza przedziału, widzimy, że zbieżność obu części nie zależy od wyboru $c$ lub $d$, a ich suma się nie zmienia.
\end{proof}
Dalej rozpatrujemy przede wszystkim prawy koniec $b\in\R\cup\{\infty\}$ przy regularnym lewym końcu $a$. Wersję dla lewego końca uzyskujemy przez odbicie $x=-t$. Dla dowolnej funkcji całkowalnej na $[u,v]$, $u<v$, odbicie punktów podziału i punktów wyróżnionych, a następnie odwrócenie ich kolejności, daje podział $[-v,-u]$ o tej samej średnicy i tych samych długościach podprzedziałów w odwrotnej kolejności. Sumy Riemanna dla $f(x)$ i $f(-t)$ są identyczne po zmianie indeksów. Ponieważ odbicie jest odwracalne, twierdzenie \ref{ma1:thm:riemann-darboux-equivalence} daje całkowalność funkcji odbitej oraz
\[
 \int_u^v f(x)\,\dd x=\int_{-v}^{-u}f(-t)\,\dd t.
\]
Przechodząc do granicy lewego końca, otrzymujemy odpowiednią granicę prawego końca funkcji odbitej. Nie potrzebujemy tu ciągłości $f$.

\section{Kryterium Cauchy'ego i całki nieujemne}
\begin{twierdzenie}[Kryterium Cauchy'ego]
\label{ma1:thm:improper-cauchy}
Niech $f$ będzie całkowalna na każdym $[a,r]$, $a<r<b\le\infty$. Całka $\int_a^b f(x)\,\dd x$ jest zbieżna wtedy i tylko wtedy, gdy dla każdego $\varepsilon>0$ istnieje $r_0\in(a,b)$ takie, że
\[
 \left|\int_u^v f(x)\,\dd x\right|<\varepsilon
 \qquad(r_0<u<v<b).
\]
\end{twierdzenie}
\begin{proof}
Oznaczmy $F(r)=\int_a^r f$. Jeżeli $F(r)\to A\in\R$, wybieramy $r_0$ tak, aby $|F(r)-A|<\varepsilon/2$ dla $r>r_0$. Z addytywności $|\int_u^v f|=|F(v)-F(u)|<\varepsilon$.

Odwrotnie, wybierzmy rosnący ciąg $r_n\to b$ w $(a,b)$: dla skończonego $b$ można przyjąć $r_n=b-(b-a)/(n+1)$, a dla $b=\infty$ --- $r_n=a+n$, $n\ge1$. Z założenia ciąg $F(r_n)$ jest ciągiem Cauchy'ego, gdyż dla dużych $n,m$ oba argumenty są większe od $r_0$, a kolejność argumentów zmienia jedynie znak różnicy. Z zupełności $\R$ ciąg ma granicę $A$. Dla dowolnego $\varepsilon>0$ stosujemy założenie z $\varepsilon/2$ i wybieramy $n$ tak, aby $r_n>r_0$ oraz $|F(r_n)-A|<\varepsilon/2$. Dla każdego $r\in(r_0,b)$ (także gdy $r=r_n$) mamy
\[
 |F(r)-A|\le |F(r)-F(r_n)|+|F(r_n)-A|<\varepsilon.
\]
Stąd $F(r)\to A$.
\end{proof}
\begin{twierdzenie}
\label{ma1:thm:improper-positive}
Jeżeli $f\ge0$ i $f$ jest całkowalna na każdym $[a,r]$, $a<r<b$, to całka niewłaściwa jest zbieżna wtedy i tylko wtedy, gdy zbiór $\{\int_a^r f:a<r<b\}$ jest ograniczony z góry. Wtedy całka jest jego supremum; w przeciwnym przypadku jest rozbieżna do $+\infty$.
\end{twierdzenie}
\begin{proof}
Funkcja $F(r)=\int_a^r f$ jest niemalejąca, ponieważ dla $u<v$ różnica $F(v)-F(u)=\int_u^v f$ jest nieujemna. Gdy zbiór wartości jest ograniczony z góry, istnieje jego supremum $A$. Dla $\varepsilon>0$ wybieramy $r_0$ z $F(r_0)>A-\varepsilon$. Dla $r>r_0$ mamy $A-\varepsilon<F(r)\le A$, więc granicą jest $A$. Gdy zbiór jest nieograniczony, dla każdej liczby $M$ istnieje $r_0$ z $F(r_0)>M$, a monotoniczność daje $F(r)>M$ dla $r>r_0$.
\end{proof}

\section{Kryteria porównawcze}
\begin{twierdzenie}
\label{ma1:thm:improper-comparison}
Niech $f,g$ będą całkowalne na każdym $[a,r]$, $a<r<b$, i niech $0\le f\le g$ na pewnym końcowym odcinku $[c,b)$, $a\le c<b$. Jeżeli $\int_a^b g$ jest zbieżna, to $\int_a^b f$ jest zbieżna. Jeżeli $\int_a^b f$ jest rozbieżna, to $\int_a^b g$ jest rozbieżna.
\end{twierdzenie}
\begin{proof}
Zmiana dolnej granicy z $a$ na $c$ odejmuje skończoną całkę $\int_a^c$. Na odcinku końcowym z porównania całek właściwych mamy $0\le\int_c^r f\le\int_c^r g$. Jeśli całka z $g$ jest zbieżna, jej całki częściowe są ograniczone na mocy twierdzenia \ref{ma1:thm:improper-positive}. To samo zachodzi dla $f$, więc stosujemy to twierdzenie. Druga teza jest kontrapozycją pierwszej.
\end{proof}
\begin{wniosek}[Kryterium porównawcze graniczne]
\label{ma1:cor:improper-limit-comparison}
Przy założeniu całkowalności $f,g$ na każdym $[a,r]$, jeśli $f\ge0$, $g>0$ na pewnym końcowym odcinku i $f(x)/g(x)\to L\in(0,\infty)$ przy $x\to b$ od lewej, to obie całki mają ten sam charakter zbieżności. Gdy $L=0$, zbieżność całki z $g$ pociąga zbieżność całki z $f$. Gdy $L=+\infty$, rozbieżność całki z $g$ pociąga rozbieżność całki z $f$.
\end{wniosek}
\begin{proof}
Dla $0<L<\infty$ na dostatecznie późnym odcinku mamy $L/2<f/g<3L/2$. Stąd $(L/2)g\le f\le(3L/2)g$. Kryterium porównawcze w obu kierunkach daje równoważność, ponieważ pomnożenie przez dodatnią stałą nie zmienia zbieżności (liniowość całki właściwej i przejście do granicy). Dla $L=0$ ostatecznie $f\le g$, a dla $L=+\infty$ ostatecznie $g\le f$. Stosujemy odpowiedni kierunek kryterium.
\end{proof}
\begin{przyklad}[Całki potęgowe]
\label{ma1:ex:improper-p-integrals}
Dla $p\in\R$ całka $\int_1^\infty x^{-p}\,\dd x$ jest zbieżna dokładnie wtedy, gdy $p>1$, natomiast $\int_0^1 x^{-p}\,\dd x$ dokładnie wtedy, gdy $p<1$.
\end{przyklad}
\begin{proof}
Dla $p\ne1$ z tablicy funkcji pierwotnych \ref{ma1:prop:primitive-table} i wzoru Newtona--Leibniza mamy
\[
 \int_1^R x^{-p}\,\dd x=\frac{R^{1-p}-1}{1-p},\qquad
 \int_\eta^1 x^{-p}\,\dd x=\frac{1-\eta^{1-p}}{1-p}.
\]
Gdy $R\to\infty$, potęga $R^{1-p}$ dąży do zera dla $p>1$ i do nieskończoności dla $p<1$. Gdy $\eta\to0^+$, potęga $\eta^{1-p}$ dąży do zera dla $p<1$ i do nieskończoności dla $p>1$. W przypadkach zbieżnych wartości wynoszą odpowiednio $1/(p-1)$ i $1/(1-p)$, a w pozostałych granice są równe $+\infty$. Dla $p=1$ całki obcięte wynoszą $\log R$ i $-\log\eta$, więc obie są rozbieżne do $+\infty$.
\end{proof}

\begin{center}
\begin{tikzpicture}[x=1cm,y=2.7cm,>=Stealth]
 \draw[->] (.7,0)--(7.1,0) node[right] {$x$};
 \draw[->] (1,-.04)--(1,1.16) node[above] {$y$};
 \fill[blue!18] (1,0)--plot[domain=1:6,samples=100] (\x,{1/(\x*\x)})--(6,0)--cycle;
 \draw[blue!70!black,thick,domain=1:6.8,samples=100] plot (\x,{1/(\x*\x)});
 \draw[orange!80!black,thick,domain=1:6.8,samples=100] plot (\x,{1/\x});
 \draw[dashed] (6,0)--(6,{1/6});
 \node[below] at (1,0) {$1$};
 \node[below] at (6,0) {$R$};
 \node[above right,orange!80!black] at (3,{1/3}) {$x^{-1}$};
 \node[below right,blue!70!black] at (2.5,.16) {$x^{-2}$};
\end{tikzpicture}

\small Rysunek: pole pod $x^{-2}$ od $1$ do $R$ wynosi $1-1/R$ i jest ograniczone niezależnie od $R$. Pole pod $x^{-1}$ wynosi $\log R$ i rośnie nieograniczenie. Sama zbieżność obu funkcji do zera nie rozstrzyga zbieżności całki.
\end{center}
\begin{przyklad}
\label{ma1:ex:improper-log-test}
Całka $\int_e^\infty \dd x/[x(\log x)^p]$ jest zbieżna dokładnie wtedy, gdy $p>1$. Podstawienie $t=\log x$ na każdym $[e,R]$ daje $\int_1^{\log R}t^{-p}\,\dd t$. Ponieważ $\log R\to\infty$, teza wynika z przykładu \ref{ma1:ex:improper-p-integrals}.
\end{przyklad}

\section{Zbieżność bezwzględna i warunkowa}
\begin{definicja}
\label{ma1:def:improper-absolute}
Całka niewłaściwa $\int_a^b f$ jest bezwzględnie zbieżna, gdy $\int_a^b |f|$ jest zbieżna. Jest warunkowo zbieżna, gdy sama jest zbieżna, lecz całka z modułu jest rozbieżna.
\end{definicja}
\begin{twierdzenie}
\label{ma1:thm:improper-absolute-convergence}
Zbieżność bezwzględna pociąga zbieżność. W tym przypadku
\[
 \left|\int_a^b f(x)\,\dd x\right|\le\int_a^b |f(x)|\,\dd x.
\]
\end{twierdzenie}
\begin{proof}
Moduł jest całkowalny na każdym przedziale obciętym z twierdzenia \ref{ma1:thm:riemann-algebra}. Jeżeli całka z modułu jest zbieżna, to dla dostatecznie bliskich końcowi $b$ argumentów $u<v<b$ kryterium Cauchy'ego daje $\int_u^v |f|<\varepsilon$. Ponieważ $|\int_u^v f|\le\int_u^v |f|$, stosujemy to samo kryterium do $f$. Nierówność dla całek niewłaściwych otrzymujemy przez przejście do granicy w nierówności dla całek obciętych. Przy kilku osobliwościach stosujemy argument osobno do każdej części i następnie nierówność trójkąta.
\end{proof}
Jeśli całki z $f$ i $g$ są zbieżne, to dla $\alpha,\beta\in\R$ całka z $\alpha f+\beta g$ jest zbieżna i
\[
 \int_a^b(\alpha f+\beta g)\,\dd x
 =\alpha\int_a^b f\,\dd x+\beta\int_a^b g\,\dd x.
\]
Istotnie, na każdym wspólnym obcięciu zachodzi liniowość całek właściwych, a wszystkie granice są skończone. Przy różnych skończonych zbiorach osobliwości rozdzielamy w ich sumie; dodatkowe regularne punkty nie zmieniają wartości całek. Addytywność wynika z odpowiedniej równości przed przejściem do granicy. Nie odejmujemy dwóch rozbieżnych całek jak dwóch liczb.

\section{Kryteria Dirichleta i Abela}
\begin{twierdzenie}[Kryterium Dirichleta]
\label{ma1:thm:improper-dirichlet}
Niech $f$ będzie ciągła na $[a,b)$, a funkcja $F(r)=\int_a^r f$ ograniczona: $|F(r)|\le M$. Niech $g$ będzie klasy $C^1$ na $[a,b)$ (z jednostronną pochodną w $a$), monotoniczna i niech $g(r)\to0$ przy końcu $b$. Wtedy $\int_a^b fg$ jest zbieżna.
\end{twierdzenie}
\begin{proof}
Z podstawowego twierdzenia rachunku całkowego $F'=f$ we wnętrzu. Na każdym $[u,v]\subset(a,b)$ całkowanie przez części daje
\[
 \int_u^v fg\,\dd x=F(v)g(v)-F(u)g(u)-\int_u^v Fg'\,\dd x.
\]
Funkcja monotoniczna dążąca do zera jest albo nieujemna i nierosnąca, albo niedodatnia i niemalejąca: dla $u<v$ w pierwszym przypadku $g(u)\ge g(v)$, więc przejście do granicy $v\to b$ daje $g(u)\ge0$; w drugim przypadku analogicznie $g(u)\le0$. Jej pochodna ma odpowiednio znak niedodatni lub nieujemny, co wynika z ilorazów różnicowych. Zatem
\[
 \int_u^v |g'|\,\dd x=|g(v)-g(u)|=|g(u)|-|g(v)|.
\]
Pierwsza równość wynika z ustalonego znaku $g'$ i wzoru Newtona--Leibniza, druga z monotoniczności ku zeru. Stąd
\[
 \left|\int_u^v fg\,\dd x\right|
 \le M(|g(u)|+|g(v)|)+M(|g(u)|-|g(v)|)
 =2M|g(u)|.
\]
Ustalmy $\varepsilon>0$. Gdy $M=0$, prawa strona jest zerowa i warunek Cauchy'ego zachodzi dla każdego $r_0\in(a,b)$. Gdy $M>0$, wybieramy $r_0$ tak, aby $|g(u)|<\varepsilon/(2M)$ dla $u>r_0$. Kryterium Cauchy'ego daje zbieżność.
\end{proof}
\begin{twierdzenie}[Kryterium Abela]
\label{ma1:thm:improper-abel}
Niech $f$ będzie ciągła na $[a,b)$ i niech $\int_a^b f$ będzie zbieżna. Jeżeli $g$ jest ograniczona, monotoniczna i klasy $C^1$ na $[a,b)$, to $\int_a^b fg$ jest zbieżna.
\end{twierdzenie}
\begin{proof}
Ograniczona funkcja monotoniczna ma skończoną granicę $L$ przy końcu $b$: dla niemalejącej jest nią supremum wartości, dla nierosnącej infimum (argument jak w twierdzeniu \ref{ma1:thm:improper-positive}). Funkcja $h=g-L$ jest monotoniczna, klasy $C^1$ i dąży do zera. Funkcja $F(r)=\int_a^r f$ jest ograniczona: z istnienia skończonej granicy jest ograniczona na końcowym odcinku, a na pozostałym domkniętym przedziale jest ciągła z twierdzenia \ref{ma1:thm:riemann-fundamental}, więc ograniczona. Kryterium Dirichleta stosowane do $f,h$ daje zbieżność $\int_a^b fh$. Ponieważ $fg=fh+Lf$, liniowość zbieżnych całek daje tezę.
\end{proof}
Są to wersje kryteriów dla ciągłej $f$ i monotonicznej $g$ klasy $C^1$, nie ogólniejsze wersje dla funkcji o wahaniu ograniczonym.

\begin{przyklad}[Zbieżność warunkowa]
\label{ma1:ex:improper-sinx-over-x}
Całka $\int_1^\infty \sin x/x\,\dd x$ jest zbieżna warunkowo.
\end{przyklad}
\begin{proof}
Dla $f(x)=\sin x$ mamy $\int_1^r f=\cos1-\cos r$, więc $|F(r)|\le2$. Funkcja $g(x)=1/x$ maleje do zera i jest klasy $C^1$, więc kryterium Dirichleta daje zbieżność. Dla każdego całkowitego $n\ge1$ i $x\in[n\pi,(n+1)\pi]$ mamy $1/x\ge1/((n+1)\pi)$. Podstawiając $x=t+n\pi$ i korzystając ze wzorów dodawania, otrzymujemy
\[
 \int_{n\pi}^{(n+1)\pi}|\sin x|\,\dd x
 =\int_0^\pi\sin t\,\dd t=2.
\]
Wobec tego
\[
 \int_\pi^{(N+1)\pi}\frac{|\sin x|}{x}\,\dd x
 \ge\frac2\pi\sum_{n=1}^{N}\frac1{n+1}.
\]
Suma po prawej rośnie nieograniczenie: grupując składniki szeregu harmonicznego o indeksach od $2^j$ do $2^{j+1}-1$, każda grupa ma sumę co najmniej $1/2$. Z twierdzenia \ref{ma1:thm:improper-positive} wynika rozbieżność całki z modułu. Dodanie zwykłej całki na $[1,\pi]$ nie zmienia wniosku.
\end{proof}

\section{Wartość główna Cauchy'ego}
\begin{definicja}
\label{ma1:def:cauchy-principal-value}
Dla funkcji całkowalnej na wszystkich domkniętych przedziałach w $\R$ wartością główną Cauchy'ego na całej prostej nazywamy skończoną granicę
\[
 \operatorname{PV}\int_{-\infty}^{\infty}f(x)\,\dd x
 :=\lim_{R\to\infty}\int_{-R}^R f(x)\,\dd x,
\]
jeżeli istnieje. Dla pojedynczej osobliwości $c\in(a,b)$ przy skończonych regularnych końcach, zakładając całkowalność na każdym $[a,c-\eta]$ i $[c+\eta,b]$ dla $0<\eta<\min(c-a,b-c)$, definiujemy
\[
 \operatorname{PV}\int_a^b f(x)\,\dd x
 :=\lim_{\eta\to0^+}\left(
 \int_a^{c-\eta}f(x)\,\dd x+\int_{c+\eta}^b f(x)\,\dd x\right),
\]
o ile granica istnieje i jest skończona.
\end{definicja}
W obu sytuacjach objętych definicją, jeżeli zwykła całka niewłaściwa jest zbieżna, to jej wartość główna istnieje i jest jej równa, ponieważ obie części mają skończone, niezależne granice. Odwrotna implikacja jest fałszywa.
\begin{przyklad}
\label{ma1:ex:principal-value-counterexamples}
Na $[-1,1]$ funkcja $1/x$ poza zerem ma całki obcięte
\[
 \int_{-1}^{-\eta}\frac{\dd x}{x}=\log\eta,\qquad
 \int_\eta^1\frac{\dd x}{x}=-\log\eta.
\]
Ich suma jest zerowa dla każdego $0<\eta<1$, więc wartość główna wynosi $0$. Każda część jest jednak rozbieżna. Jeżeli odcinamy niezależnie odcinki długości $\delta$ i $\eta$ z lewej i prawej strony zera, suma wynosi $\log(\delta/\eta)$. Przy $\delta=k\eta$ dla $k>0$ i dostatecznie małym $\eta$ (tak aby $\delta<1$) granicą jest $\log k$. Ponieważ dla dowolnego $A\in\R$ można wziąć $k=e^A$, przez taki dobór obcięć otrzymujemy dowolną wartość rzeczywistą. Nie jest to niezależna granica obu części.

Podobnie $\operatorname{PV}\int_{-\infty}^{\infty}x/(1+x^2)\,\dd x=0$ ze względu na nieparzystość (podstawienie $x=-t$). Zwykła całka jest rozbieżna: całka od $0$ do $R$ wynosi $\tfrac12\log(1+R^2)\to+\infty$, a część lewa ma granicę $-\infty$.
\end{przyklad}
Wartość główna jest osobną konstrukcją z określoną zależnością parametrów obcięcia. Jej istnienie nie oznacza zbieżności warunkowej.

\section{Podstawienie i całkowanie przez części}
\begin{twierdzenie}
\label{ma1:thm:improper-substitution}
Niech $f$ będzie ciągła na $[a,b)$, gdzie $b$ może być nieskończony. Niech $\alpha\in\R$, $\alpha<\beta\le\infty$ oraz $\varphi:[\alpha,\beta)\to[a,b)$ będzie klasy $C^1$, niemalejąca, $\varphi(\alpha)=a$ oraz $\varphi(t)\to b$ przy $t\to\beta$ od lewej, gdzie dopuszczamy $\beta=\infty$. Wtedy całki
\[
 \int_a^b f(x)\,\dd x,\qquad
 \int_\alpha^\beta f(\varphi(t))\varphi'(t)\,\dd t
\]
są jednocześnie zbieżne i w przypadku zbieżności równe.
\end{twierdzenie}
\begin{proof}
Przedłużamy $f$ ciągle na otwarty przedział $(a-1,b)$, przyjmując $f(x)=f(a)$ dla $x<a$. Dla $\alpha<T<\beta$ funkcja $\varphi$ jest klasy $C^1$ na $[\alpha,T]$, a jej wartości należą do tego otwartego przedziału. Wzór podstawienia dla całek właściwych (twierdzenie \ref{ma1:thm:integration-substitution}) daje
\[
 \int_\alpha^T f(\varphi(t))\varphi'(t)\,\dd t
 =\int_a^{\varphi(T)} f(x)\,\dd x.
\]
Równość obejmuje również przypadek $\varphi(T)=a$, w którym prawa strona jest zerowa. Jeżeli całka z $f$ jest zbieżna, przechodzimy do granicy $T\to\beta$. Odwrotnie, niech lewa strona ma granicę $A$. Ustalmy $\varepsilon>0$ i $T_0$ tak, aby dla $T>T_0$ lewa strona różniła się od $A$ o mniej niż $\varepsilon$. Dla każdego $r\in(\varphi(T_0),b)$ istnieje $T_1>T_0$ z $\varphi(T_1)>r$, ponieważ $\varphi(T)\to b$. Z ciągłości istnieje $T\in(T_0,T_1)$ z $\varphi(T)=r$. Równość całek obciętych daje $|\int_a^r f-A|<\varepsilon$. To dowodzi zbieżności całki z $f$.
\end{proof}
Monotoniczność jest wystarczająca, lecz w tym dowodzie nie jest niezbędna: wystarczą ciągłość, klasa $C^1$ na obcięciach, wskazany obraz i warunki graniczne. Wersję dla niewłaściwego lewego końca uzyskujemy przez odbicie, a podstawienie malejące z odpowiednio odwróconymi końcami wymaga zmiany orientacji. Podstawienie nie może łączyć dwóch niezależnych granic przy osobliwościach w jedną granicę bez uzasadnienia.
\begin{twierdzenie}
\label{ma1:thm:improper-parts}
Niech $u,v$ będą klasy $C^1$ na $[a,b)$. Jeżeli istnieje skończona granica $\ell=\lim_{r\to b^-}u(r)v(r)$ (przy $b=\infty$ rozumiemy $r\to\infty$) oraz całka $\int_a^b u'v$ jest zbieżna, to $\int_a^b uv'$ jest zbieżna i
\[
 \int_a^b uv'\,\dd x=\ell-u(a)v(a)-\int_a^b u'v\,\dd x.
\]
Jeżeli granica iloczynu istnieje i jest skończona, zbieżność którejkolwiek całki pociąga zbieżność drugiej.
\end{twierdzenie}
\begin{proof}
Dla $a<r<b$ całkujemy przez części na $[a,r]$:
\[
 \int_a^r uv'+\int_a^r u'v=u(r)v(r)-u(a)v(a).
\]
Przenosimy całkę, o której zakładamy zbieżność, na prawą stronę i przechodzimy do granicy. Wszystkie składniki prawej strony mają skończone granice, więc istnieje druga całka i zachodzi wskazana równość. Argument jest symetryczny.
\end{proof}
\begin{przyklad}
\label{ma1:ex:improper-parts-exp}
Obliczymy $\int_0^\infty xe^{-x}\,\dd x$. Dla $u=x$, $v=-e^{-x}$ mamy $u'=1$, $v'=e^{-x}$. Z szeregu wykładniczego dla $R>0$ mamy $e^R\ge R^2/2$, więc $0\le Re^{-R}\le2/R\to0$. Ponadto $\int_0^R e^{-x}\,\dd x=1-e^{-R}\to1$. Wobec twierdzenia \ref{ma1:thm:improper-parts}
\[
 \int_0^\infty xe^{-x}\,\dd x=0-0+\int_0^\infty e^{-x}\,\dd x=1.
\]
\end{przyklad}
Przy badaniu całki niewłaściwej najpierw ustalamy wszystkie niewłaściwe końce i punkty osobliwe. Następnie badamy każdą część niezależnie. Dopiero po ustaleniu zbieżności można sumować wartości i stosować działania na skończonych granicach. Dla funkcji nieujemnych najczęściej wystarczają porównania, a dla oscylujących istotne jest rozróżnienie zbieżności bezwzględnej i warunkowej. W następnym rozdziale przejdziemy do całek zależnych od parametru i zamiany granicy z całką.
